%% ma: L12-B11-0017
%% lop: 12
%% chuong: 4
%% bai: 11
%% dang: 12.11.11
%% loai: TN4
%% muc: TH
%% dap_an: "D"
%% nguon: "(NBV)-10 ĐỀ KIỂM TRA HK2-TOÁN 12-DUNG CHUNG 3 SÁCH.pdf (H02, Đề số 2), Phần 1 Câu 6"
%% kien_thuc: "Bài 11 (nguyên hàm của hàm lũy thừa)"
%% trang_thai: chinh-thuc
\begin{ex}
Hàm số nào sau đây \textbf{không} là một nguyên hàm của $f(x)=\sqrt[3]{x}$ trên $(0;+\infty)$?
\choice{$F_1(x)=\dfrac{3\sqrt[3]{x^4}}{4}+1$}{$F_3(x)=\dfrac{3x\sqrt[3]{x}}{4}+3$}{$F_4(x)=\dfrac34x^{\frac43}+4$}{\True $F_2(x)=\dfrac{3\sqrt[4]{x^3}}{4}+2$}
\loigiai{Ta có $f(x)=x^{\frac13}$ nên $\displaystyle\int f(x)\,\mathrm dx=\dfrac{3x^{\frac43}}{4}+C=\dfrac{3\sqrt[3]{x^4}}{4}+C=\dfrac{3x\sqrt[3]{x}}{4}+C$. Các hàm $F_1$, $F_3$, $F_4$ có dạng này. Hàm $F_2(x)=\dfrac34x^{\frac34}+2$ có đạo hàm $\dfrac9{16}x^{-\frac14}\ne f(x)$ nên không là nguyên hàm. Chọn D.}
\end{ex}
